% =====================================================================
%  3F3 Module Experiment: Random Variables and Random Number Generation
%  Lab report template
% ---------------------------------------------------------------------
%  How to use it:
%  * Fill in your details below.
%  * Replace each blue \answer{...} placeholder with your answer. The
%    grey \task{...} text summarises what each part asks for; change
%    \showpromptstrue to \showpromptsfalse to hide all of it.
%  * Figures: save each plot in the figures/ folder under the name shown
%    in its placeholder box (names are case sensitive), e.g.
%        plt.savefig("figures/q1-histograms.pdf", bbox_inches="tight")
%    It then appears automatically (.pdf, .png or .jpg). On Overleaf,
%    select the figures folder before clicking Upload.
%  * Code: paste into the lstlisting blocks, or include a file with
%        \lstinputlisting{q1.py}
%  * Overleaf: New Project > Upload Project, then choose the .zip file.
% =====================================================================
\documentclass[11pt,a4paper]{article}
\usepackage[margin=2.5cm]{geometry}
\usepackage[T1]{fontenc}
\usepackage{lmodern}
\usepackage{amsmath,amssymb}
\usepackage{graphicx}
\usepackage{xcolor}
\usepackage{listings}
\usepackage{parskip}
\usepackage[section]{placeins}   % keeps figures within their question

% ---- Your details ---------------------------------------------------
\newcommand{\studentname}{\answer{Your name}}
\newcommand{\crsid}{\answer{CRSID}}
\newcommand{\college}{\answer{College}}
\newcommand{\labdate}{\answer{Date of lab session}}

\newif\ifshowprompts
\showpromptstrue   % \showpromptsfalse hides the grey task prompts

% ---- Template helpers (no need to edit) -----------------------------
\definecolor{answerblue}{rgb}{0.1,0.25,0.7}
\definecolor{taskgrey}{gray}{0.4}

% \task{...}: what the handout asks for in this part
\NewDocumentCommand{\task}{+m}{%
  \ifshowprompts{\small\itshape\leavevmode\color{taskgrey}#1\par}\fi}

% \answer{...}: placeholder, to be replaced by your answer
\NewDocumentCommand{\answer}{+m}{\leavevmode{\color{answerblue}[#1]}}

% \answerfig[options]{name}{caption}: shows figures/name.pdf (or .png,
% .jpg, .jpeg) if it exists, otherwise a placeholder box. Options are
% passed to \includegraphics. Refer to it as Figure~\ref{fig:name}.
\NewDocumentCommand{\answerfig}{O{width=0.7\textwidth} m m}{%
  \begin{figure}[htbp]
    \centering
    \IfFileExists{figures/#2.pdf}{\includegraphics[#1]{figures/#2.pdf}}{%
    \IfFileExists{figures/#2.png}{\includegraphics[#1]{figures/#2.png}}{%
    \IfFileExists{figures/#2.jpg}{\includegraphics[#1]{figures/#2.jpg}}{%
    \IfFileExists{figures/#2.jpeg}{\includegraphics[#1]{figures/#2.jpeg}}{%
      \fbox{\parbox[c][0.22\textwidth][c]{0.7\textwidth}{\centering
        \color{answerblue}Put your plot in the \texttt{figures} folder as\\
        \texttt{\detokenize{#2}.pdf} or \texttt{\detokenize{#2}.png}}}}}}}%
    \caption{#3}\label{fig:#2}%
  \end{figure}}

% Unicode symbols typed straight into the text
\DeclareUnicodeCharacter{03BC}{\ensuremath{\mu}}
\DeclareUnicodeCharacter{03C3}{\ensuremath{\sigma}}
\DeclareUnicodeCharacter{2212}{\ensuremath{-}}
\DeclareUnicodeCharacter{2248}{\ensuremath{\approx}}
\DeclareUnicodeCharacter{2264}{\ensuremath{\leq}}
\DeclareUnicodeCharacter{2265}{\ensuremath{\geq}}

\lstset{
  language=Python,
  basicstyle=\small\ttfamily,
  keywordstyle=\color{blue!60!black},
  commentstyle=\color{green!40!black},
  stringstyle=\color{orange!60!black},
  showstringspaces=false,
  frame=single,
  framerule=0.3pt,
  breaklines=true,
  columns=fullflexible,
  keepspaces=true,
  upquote=true,
  tabsize=4,
  % non-ASCII characters in pasted code
  literate={μ}{{$\mu$}}1 {σ}{{$\sigma$}}1 {π}{{$\pi$}}1 {²}{{$^2$}}1
           {×}{{$\times$}}1 {≈}{{$\approx$}}1 {≤}{{$\leq$}}1 {≥}{{$\geq$}}1
           {−}{{-}}1 {–}{{-}}1 {—}{{--}}1
           {‘}{{\textquotesingle}}1 {’}{{\textquotesingle}}1
           {“}{{"}}1 {”}{{"}}1,
}

\begin{document}

\begin{center}
  \textbf{Engineering Tripos Part IIA \qquad Module Experiment 3F3}\\[1.5ex]
  {\Large\bfseries Random Variables and Random Number Generation}\\[1ex]
  {\large Lab Report}\\[3ex]
  \begin{tabular}{@{}ll@{}}
    Name:        & \studentname \\
    CRSID:       & \crsid \\
    College:     & \college \\
    Lab session: & \labdate \\
  \end{tabular}
\end{center}
\medskip
\hrule

% =====================================================================
\section{Uniform random variables}
% =====================================================================

\subsection*{Five numbers from seed 1234}

\task{Initialise an \texttt{LCG} random number generator with seed $1234$
and print five uniform random numbers. What numbers do you get?}

\begin{lstlisting}
  generator = LCG(1234)

  generator.generate()
  generator.generate()
  generator.generate()
  generator.generate()
  generator.generate()

\end{lstlisting}

\answer{0.7143076679203659
  0.20701311994343996
  0.7495418272446841
  0.3460624306462705
  0.8034394562710077. The numbers perhaps don't seem particularly random. Alternate numbers are increasing from 0.71 by about 0.05 at a time, while the other pattern is increasing from 0.2. This is because of the relative size of the seed to the size of the modulo.}

\subsection*{Histograms}

\task{Generate $1{,}000$ uniform random variables and plot a histogram, then
do the same for $100{,}000$. Provide detailed comments on what you observe.}

\begin{lstlisting}
# generate uniform numbers
X1 = np.array([generator.uniform() for _ in range(1000)])
X2 = np.array([generator.uniform() for _ in range(100000)])
# Plot histograms
fig, ax = plt.subplots(2)
ax[0].hist(X1, bins=np.linspace(0,1,20))
ax[1].hist(X2, bins=np.linspace(0,1,20))
plt.show()
\end{lstlisting}

\answerfig[width=\textwidth]{q1-histograms}{Histograms of $1{,}000$ and
$100{,}000$ uniform random numbers generated by the LCG.}

\answer{With 1000 samples, this does not well model a uniform distribution. The deterministic nature of the process is evident from the fact that the samples cluster around certain buckets. There are 4 or 5 distinct intervals where there are 50-60 samples, while the rest of the buckets are closer to 40. We can also say that the entropy is lower than a uniform distribution as a sample is significantly more likely to be in one of the clusters than elsewhere in the distribution. For the 100,000 samples case, the distribution is now much closer to uniform. However, there are still some undulations visible, that line up with the same intervals that peaked with 1000 samples. The undulations have much lower magnitudes and are much more gradual.}

\subsection*{Mean, variance and covariance}
\answer{Theoretically, $\mathbb{E}[\mathcal{U}(0,1)] = \int_0^1xf(x)dx=\int_0^1xdx=[\frac{1}{2}x^2]_0^1=\frac{1}{2}$

For the variance, $\mathrm{Var}(X)=\mathbb{E}[X^2]-\mathbb{E}[X]^2$. First,
$\mathbb{E}[X^2]=\int_0^1x^2f(x)dx=\int_0^1x^2dx=[\frac{1}{3}x^3]_0^1=\frac{1}{3}$.
Therefore $\mathrm{Var}(X)=\frac{1}{3}-\left(\frac{1}{2}\right)^2=\frac{1}{3}-\frac{1}{4}=\frac{1}{12}\approx 0.0833$.}\\
\task{Calculate the mean and variance of the $100{,}000$ numbers and compare
them with the theoretical values for $\mathcal{U}(0,1)$. Compute the
covariance between $X_{t+1}$ and $X_t$. Why does this provide evidence that
the numbers are behaving like uniform random numbers?}

\begin{lstlisting}
#calculate mean
mean_x2 = np.mean(X2)
print(f"X2 mean = {mean_x2}")

#calculate variance
variance_x2 = np.var(X2)
print(f"X2 variance = {variance_x2}")

#calculate covariance
X_t = X2[:-1] #starts at the start
X_t1 = X2[1:] #starts offset by one
covariance = np.cov(X_t, X_t1)[0,1] #covariance matrix, this index gives the Cov.
print(f"X2 covariance = {covariance}")
\end{lstlisting}

\begin{center}
  \begin{tabular}{lcc}
    \hline
             & LCG, $100{,}000$ samples & Theory, $\mathcal{U}(0,1)$ \\
    \hline
    Mean     & \answer{0.499} & $\frac{1}{2}=0.5$ \\
    Variance & \answer{0.083} & $\frac{1}{12}\approx 0.0833$ \\
    \hline
  \end{tabular}
\end{center}

Covariance between $X_{t+1}$ and $X_t$: \answer{0.0001}

\answer{The mean and variance are very close to the theoretical values for the uniform distribution. These are evidence that the numbers follow a uniform distribution but they do not prove it. The extremely low covariance is stronger evidence. A uniform distribution is memoryless, and the low covariance shows that the next observation is independent of previous observations, satisfying the memoryless property.}

% =====================================================================
\section{Functions of random variables}
% =====================================================================

\subsection*{\boldmath $y = x^2$}

\task{For $x \sim \mathcal{U}(0,1)$ and $y = g(x) = x^2$, calculate $p_Y(y)$
using the Jacobian formula. Verify it by transforming a large collection of
samples $x^{(i)}$ to give $y^{(i)} = g(x^{(i)})$, histogramming the $y$
samples and overlaying your formula.}

\answer{
The pdf of $x \sim \mathcal{U}(0,1)$ is
\[
p_X(x) = \begin{cases} 1 & 0 < x < 1 \\ 0 & \text{otherwise,} \end{cases}
\]
so the support of $x$ is $(0,1)$. Differentiating $g(x) = x^2$ gives
\[
\frac{dg}{dx} = 2x, \qquad \left|\frac{dg}{dx}\right| = 2|x| ,
\]
which is nonzero everywhere on the support of $x$.

For a given $y > 0$, the equation $y = x^2$ has two inverses,
\[
x_1(y) = \sqrt{y}, \qquad x_2(y) = -\sqrt{y}.
\]
Hence, by the Jacobian formula, $p_Y(y)$ is a sum over both inverses,
\[
p_Y(y) = \sum_{k=1}^{2} \frac{p_X\!\left(x_k(y)\right)}{\left|\frac{dg}{dx}\right|_{x=x_k(y)}}
= \frac{p_X(\sqrt{y})}{2\sqrt{y}} + \frac{p_X(-\sqrt{y})}{2\sqrt{y}} ,
\]
since $|2x_k| = 2\sqrt{y}$ for both $k$.

Now evaluate $p_X$ at each inverse. For $x_1 = \sqrt{y}$, $p_X(\sqrt{y}) = 1$ when $0 < y < 1$, and $0$ otherwise. For $x_2 = -\sqrt{y}$, we have $-\sqrt{y} < 0$ for all $y > 0$, which lies outside the support of $x$, so $p_X(-\sqrt{y}) = 0$. Only the $k=1$ term contributes, giving
\[
p_Y(y) = \begin{cases} \dfrac{1}{2\sqrt{y}} & 0 < y < 1 \\[1ex] 0 & \text{otherwise.} \end{cases}
\]
As a check, $\int_0^1 \frac{1}{2\sqrt{y}}\,dy = \left[\sqrt{y}\right]_0^1 = 1$, so $p_Y$ is a valid pdf.
}

\begin{lstlisting}
# transform the LCG samples
X2 = np.array([generator.uniform() for _ in range(100000)])
Y2 = X2 ** 2

y_grid = np.linspace(0.001, 1, 1000)
p_Y = 1 / (2 * np.sqrt(y_grid)) #per derivation

# build the chart
fig, ax = plt.subplots()
ax.hist(Y2, bins=np.linspace(0, 1, 50), density=True, label="Histogram of $y=x^2$ samples")
ax.plot(y_grid, p_Y, 'r', label=r"$p_Y(y)=\frac{1}{2\sqrt{y}}$")
ax.set_xlabel("y")
ax.set_ylabel("$p_Y(y)$")
ax.set_ylim(0, 5)  # p_Y diverges at y=0
ax.legend()
fig.savefig("figures/q2-square.pdf")
plt.show()
\end{lstlisting}

\answerfig{q2-square}{Histogram of $y = x^2$ for transformed uniform samples,
with the density from the Jacobian formula overlaid.}

\answer{This is quite a close fit to the theoretical distribution. Even towards the asymptote, it looks very close, but I would expect it to not follow the curve very high up as the number of samples and bin width are finite, while the curve tends to infinity.}

\subsection*{\boldmath $y = \ln(x) - \ln(1-x)$}

\task{Calculate $p_Y(y)$ for $y = g(x) = \ln(x) - \ln(1-x)$ using the
Jacobian formula, and again verify your result by histogramming transformed
samples.}

\answer{
As before, the pdf of $x \sim \mathcal{U}(0,1)$ is
\[
p_X(x) = \begin{cases} 1 & 0 < x < 1 \\ 0 & \text{otherwise,} \end{cases}
\]
so the support of $x$ is $(0,1)$. Differentiating $g(x) = \ln(x) - \ln(1-x)$ gives
\[
\frac{dg}{dx} = \frac{1}{x} + \frac{1}{1-x} = \frac{1}{x(1-x)}, \qquad \left|\frac{dg}{dx}\right| = \frac{1}{x(1-x)},
\]
which is positive and finite everywhere on the support of $x$. Hence $g$ is strictly increasing on $(0,1)$ and has exactly one inverse for each $y$.

$y = \ln\!\left(\frac{x}{1-x}\right)$, so $e^{y} = \frac{x}{1-x}$ and
\[
x(y) = \frac{e^{y}}{1+e^{y}} = \frac{1}{1+e^{-y}}.
\]
As $x \to 0^+$, $y \to -\infty$, and as $x \to 1^-$, $y \to +\infty$, so $x_1(y) \in (0,1)$ for every $y \in \mathbb{R}$. Hence $p_X(x(y)) = 1$ for all $y$.

By the Jacobian formula, $p_Y(y)$ is:
\[
p_Y(y) = \frac{p_X\!\left(x(y)\right)}{|\frac{dg}{dx}|} = x(1-x).
\]
Using $x_1 = \frac{e^y}{1+e^y}$ and $1-x_1 = \frac{1}{1+e^y}$,
\[
p_Y(y) = \frac{e^{y}}{\left(1+e^{y}\right)^2}, \qquad y \in \mathbb{R}.
\]
The support of $y$ spans all real numbers which is the image of $(0,1)$ under $g$. Checking as before, $\int_{-\infty}^{\infty} p_Y(y)\,dy = \left[\frac{e^y}{1+e^y}\right]_{-\infty}^{\infty} = 1 - 0 = 1$, so $p_Y$ is a valid pdf.
}

\begin{lstlisting}
    # transform fresh LCG samples
    X2 = np.array([generator.uniform() for _ in range(100000)])
    Y3 = np.log(X2) - np.log(1 - X2)

    y_grid = np.linspace(-10, 10, 1000)
    p_Y = np.exp(y_grid) / (1 + np.exp(y_grid)) ** 2 #per derivation

    # build the chart
    fig, ax = plt.subplots()
    ax.hist(Y3, bins=np.linspace(-10, 10, 80), density=True, label=r"Histogram of $y=\ln(x)-\ln(1-x)$ samples")
    ax.plot(y_grid, p_Y, 'r', label=r"$p_Y(y)=\frac{e^y}{(1+e^y)^2}$")
    ax.set_xlabel("y")
    ax.set_ylabel("$p_Y(y)$")
    ax.legend()
    fig.savefig("figures/q2-log.pdf")     
    plt.show()
\end{lstlisting}

\answerfig{q2-log}{Histogram of $y = \ln(x) - \ln(1-x)$ for transformed
uniform samples, with the density from the Jacobian formula overlaid.}

\answer{The histogram closely resembles the theoretical distribution, which is the logistic distribution (looks similar to normal, but with a sharper peak). There is one bucket that is slightly off (to the right of the peak) but that can easily be explained by the fact that the initial "uniform" distribution is not actually uniform.}

% =====================================================================
\section{Non-uniform random number generation}
% =====================================================================

\subsection*{Why the inverse CDF method works}

\task{To generate random variables with CDF $F(y)$, we generate
$x \sim \mathcal{U}(0,1)$ and set $y = F^{-1}(x)$. Why does this work?}

\answer{
The method works because the transformation $y = F^{-1}(x)$ turns a uniform density into the density whose CDF is $F$. We can show this with the Jacobian formula from the previous part. $F$, being a CDF, is strictly increasing over the support of $Y$, so $F^{-1}$ exists. The inverse of $y = F^{-1}(x)$ is then simply $x = F(y)$. Differentiating, $\frac{dx}{dy} = p_Y(y)$, since the pdf is the derivative of the CDF. Since $x \sim \mathcal{U}(0,1)$ has $p_X(x) = 1$,
\[
p_Y(y) = \frac{p_X\!\left(F(y)\right)}{\left|\frac{dy}{dx}\right|} = \frac{1}{1/p_Y(y)} = p_Y(y).
\]
So this gives exactly the density we wanted.
}

\subsection*{CDF and inverse CDF of the exponential distribution}

\task{For the exponential distribution with mean one,
$p_Y(y) = \exp(-y)$, $y \geq 0$, calculate the CDF and the inverse CDF.}

\answer{
For $y \ge 0$, the CDF is
\[
F(y) = \int_0^{y} e^{-t}\,dt = \left[-e^{-t}\right]_0^{y} = 1 - e^{-y}, \qquad 0 \le y < 1
\]
and $F(y) = 0$ for $y < 0$. To find the inverse, set $x = 1 - e^{-y}$, so $e^{-y} = 1 - x$ and
\[
F^{-1}(x) = -\ln(1-x), \qquad 0 \le x < 1.
\]
}

\subsection*{Sampling from the exponential distribution}

\task{Use the inverse CDF method to generate samples from the exponential
distribution. Plot histograms or kernel density estimates and overlay them on
the exponential density, and comment on the results.}

\begin{lstlisting}
    X4 = np.array([generator.uniform() for _ in range(100000)])
    Y4 = -np.log(1 - X2)

    y_grid = np.linspace(0, 8, 1000)
    p_Y = np.exp(-y_grid) #exponential pdf with mean 1

    # build the chart
    fig, ax = plt.subplots()
    ax.hist(Y4, bins=np.linspace(0, 8, 60), density=True, label=r"Histogram of $y=-\ln(1-x)$ samples")
    ax.plot(y_grid, p_Y, 'r', label=r"$p_Y(y)=e^{-y}$")
    ax.set_xlabel("y")
    ax.set_ylabel("$p_Y(y)$")
    ax.legend()
    fig.savefig("figures/q3-exponential.pdf")
    plt.show()
\end{lstlisting}

\answerfig{q3-exponential}{Samples generated by the inverse CDF method,
compared with the exponential density $p_Y(y) = \exp(-y)$.}

\answer{The close relationship between the histogram and the derived distribution is evidence that this method works.}

% =====================================================================
\section{Simulating a stochastic process}
% =====================================================================

\subsection*{Bernoulli and binomial random variables}

\task{Complete the \texttt{bernoulli} method of the \texttt{LCG} class to
generate Bernoulli random variables with parameter $p$, and add a method that
generates binomial random variables.}

\answer{See the code comments for explanation.}

\begin{lstlisting}
import numpy as np
import matplotlib.pyplot as plt

class LCG(object):
    def __init__(self, seed):
        self.X = int(seed)

    def generate(self):
        self.X = (1664525 * self.X + 1013904223) % 4294967296
        return self.X
    
    def uniform(self):
        return self.generate() / 4294967296
    
    def bernoulli(self, p):
        # Draw from a uniform (0,1) RV. Return 1 is the result is below p or 0 if the result is above p.
        # Because we are drawing from a uniform distribution, the probability that the result is less than p is p.
        # And therefore the probability of drawing above p is 1-p. So this acts just like a bernoulli RV.
        draw = self.uniform()

        if draw < p:
            return 1
        
        return 0

    def binomial(self, n, p):
        # Generate n independent bernoulli trials and sum them together
        trials = [self.bernoulli(p) for _ in range(n)]
        return sum(trials)
\end{lstlisting}

\subsection*{Neutral selection}

\task{Write Python code to simulate
$N_A(t+1) \sim \mathrm{Binomial}\big(N, N_A(t)/N\big)$,
$N_B(t+1) = N - N_A(t+1)$. Run simulations with $N = 200$ and
$N_A(0) = 100$ until $t = 500$, plot $N_A(t)$ and $N_B(t)$ over time, and
repeat this several times.}

\begin{lstlisting}
# Your code for this part
\end{lstlisting}

\answerfig[width=\textwidth]{q4-neutral-selection}{$N_A(t)$ and $N_B(t)$ for
several runs of the neutral selection simulation, with $N = 200$ and
$N_A(0) = 100$.}

\answer{Your comments on what you observe.}

\end{document}
